\documentclass[11pt,a4paper]{article}
\usepackage[T1]{fontenc}
\usepackage{lmodern,amsmath,amssymb,booktabs,microtype}
\usepackage[margin=24mm]{geometry}
\usepackage[colorlinks=true,urlcolor=blue,linkcolor=blue]{hyperref}
\newcommand{\EE}{\mathbb E}
\newcommand{\tr}{\operatorname{tr}}
\newcommand{\Rem}{\operatorname{Rem}}
\newcommand{\W}{\mathsf W_2}
\title{Coherent adjacent-charge products in the memory core\\The cross with the derivative singlet}
\author{Alan Day \and Chengheng\thanks{AI research assistant (OpenAI).}}
\date{Local research note --- 6 November 2026}
\begin{document}\maketitle
\begin{abstract}
We compute the previously missing equal-time cross between the
output-rank deficit and the derivative singlet $A$. Adjacent-charge
covariance composition gives a coherent Racah transform. Its signed
first moment has an exact contiguous reduction, which preserves the
cancellation of two terms of order $m$. At half filling the cross,
after the energy-square subtraction, equals the previously obtained
output-deficit variance. The next calculation is the norm of their
difference in this same sector.
\end{abstract}
\section{Operators and the result}
Use standard CAR, the full overlapping complex Gaussian $p$-body
ensemble, $\hbar=1$, and the normalized maximally mixed trace in
charge $m=N/2$. Inner products are $\langle U,V\rangle_m=
\EE\tr_m(U^\dagger V)$. Write $h=H_m$, $k=H_{m-1}$ for the
\emph{same} couplings, $C=\EE h^2$, and
\begin{align}
 A&=\Rem(k^2)+(2p-m)h^2,
 &\Rem(X)&=\sum_i c_i^\dagger Xc_i,\notag\\
 O&=\W A, &B&=\W h^2=h^2-CI.\label{ops}
\end{align}
Here $\W$ denotes ordinary Gaussian second chaos, so $O=A-\EE A$.
For input ranks $a\leq b$, let $Z^+_{ab;r}=P_r(h_a\odot h_b)$,
where $\odot$ is a square for $a=b$ and an anticommutator otherwise.
Set $S=a+b-r$ and $\mathcal S_N=\sum S\W Z^+_{ab;r}$.
Let $D_N=\sum_a(p-a)h_a$, $\bar s=\sum_a(p-a)C_a/C$, and
$L_N=\W\{h,D_N-\bar s h\}$.

First take $p^2/N\to\lambda>0$, with physical $C\to J_s^2/\lambda$;
put $\theta=2\lambda$, $q=e^{-4\lambda}$. We obtain
\begin{align}
 C^{-2}\langle B,O\rangle_m&\longrightarrow\theta(1-q),\notag\\
 C^{-2}\langle L_N,O\rangle_m&\longrightarrow0,\notag\\
 C^{-2}\langle\mathcal S_N,O\rangle_m
 &\longrightarrow\theta^2(1+q)+\theta(1-q).\label{main}
\end{align}
The final line is the output-deficit norm from the preceding note.
The proof below computes the actual cross; it does not substitute an
independent disorder copy for the two factors of $A$.

\newpage
\section{A coherent covariance transform}
Temporarily allow general $m$, and use independent, isotropic
Hermitian rank-$a$ and rank-$b$ tensors $U,V$. Their restrictions
$U^\pm,V^\pm$ to charges $m,m-1$ are coherently related primitive
tensors, each normalized to $\EE(U^\pm)^2=I$.
Let $R=\min(m-1,N-m)$, $D_+=\binom Nm$, $D_-=\binom N{m-1}$,
and define
\begin{align}
 d_r&=\binom Nr^2-\binom N{r-1}^2,
 &e_t&=\binom Nt\left[\binom N{t+1}-\binom N{t-1}\right],\notag\\
 \ell_r^2&=\frac{m(m-r)(N-m+1-r)}{N-m+1}.\label{dims}
\end{align}
The charged $(t+1,t)$ covariance eigenvalue of
$\Psi_a(X)=\EE U^+XU^-$ is $\zeta_a(t)=\ell_a T_a(t)/m$, where
\begin{equation}
 T_a(t)={}_4F_3\!\left(\begin{matrix}-a,a-N-1,-t,t-N\\
 1-m,m-N,-N-1\end{matrix};1\right).\label{T}
\end{equation}
One way to fix this normalization uses the known embedded-ensemble
charged covariance, $\binom{m-t-1}{p}\binom{N-m+p-t}{p}$.
If $\omega_a^{(p)}$ are the usual upper-sector rank weights, its
triangular pure-rank decomposition reads
\begin{equation}
 \sum_{a=0}^p\omega_a^{(p)}T_a(t)=
 \frac{m\binom{m-t-1}{p}\binom{N-m+p-t}{p}}
 {(m-p)\binom mp\binom{N-m+p}{p}}.\label{triangular}
\end{equation}
Substitution of the terminating series and interchange of the finite
$(a,j)$ sums gives this binomial identity (Vandermonde summation).
The coefficient of $T_p$ is nonzero, so the triangular system fixes
all pure-rank eigenvalues. This uses established embedded-ensemble
recoupling machinery, not a new ensemble prescription \cite{kota}.

Define the coherent output contraction
\begin{equation}
 M_{ab}(r)=\EE\langle P_r(U^+V^+),\Rem P_r(U^-V^-)\rangle_m,
 \qquad \mu_{ab}(r)=\frac{mM_{ab}(r)}{\ell_a\ell_b}.
\end{equation}
Composition of $\Psi_a,\Psi_b$, followed by covariance reshuffling,
gives
\begin{align}
 T_a(t)T_b(t)&=\sum_r\mu_{ab}(r)T_r(t),\notag\\
 \mu_{ab}(r)&=\frac{d_r\ell_r^2}{m^2D_+D_-}
       \sum_{t=0}^R e_tT_a(t)T_b(t)T_r(t).\label{mu}
\end{align}
Indeed reshuffling is an HS isometry and
$\sum_t e_t\zeta_r(t)\zeta_s(t)=D_+D_-\delta_{rs}/d_r$.
This proves the inversion without selecting tensor-product
multiplicity bases or discarding their phases. We only use
$\sum_r\mu_{ab}(r)=1$; no positivity assumption is needed.

\newpage
\section{Exact unsigned and signed first moments}
Return to $N=2m$, $a,b<m$. Put $x_r=r(N+1-r)$,
$X=x_{a+b}-x_r$, and let $\chi_{ab}$ be the previous neutral
Racah kernel, with lower parameters $(-m,-m,-N-1)$.
Write
\begin{equation}
 \xi_{ab}={}_4F_3\!\left(\begin{matrix}-a,a-N-1,-b,b-N-1\\
 1-m,-m,-N-1\end{matrix};1\right),\quad
 C_f=\frac{m^2(N+1)}N,\quad
 C_t=\frac{m(m-1)(N+1)}{N-1}.
\end{equation}
With $\varepsilon=(-1)^{a+b-r}$, the four moments are
\begin{align}
 \sum_r\mu_{ab}(r)&=1,&
 \sum_rX\mu_{ab}(r)&=-2ab+\frac{x_ax_b}{C_t},\notag\\
 \sum_r\varepsilon\mu_{ab}(r)&=\xi_{ab},&
 \sum_r\varepsilon X\mu_{ab}(r)
 &=-2ab\xi_{ab}+m(\xi_{ab}-\chi_{ab}).\label{moments}
\end{align}
The unsigned identity follows by evaluating \eqref{mu} at $t=1$,
since $T_a(1)=1-x_a/C_t$.

Here is a difference-equation derivation of the signed identities.
The Racah equation \cite{ks}, with
$(\alpha,\beta,\gamma,\delta)=(-m,-m-1,-N-2,0)$, yields
\begin{align}
 (\mathcal Lf)_a&=B_a(f_a-f_{a+1})+D_a(f_a-f_{a-1}),\notag\\
 B_a&=\frac{(a+1-m)(a-N-1)^2}{2(2a-N-1)},&
 D_a&=\frac{a^2(a-m-2)}{2(2a-N-1)},\label{rates}
\end{align}
with $B_{m-1}=D_0=0$. Its eigenvalue on $T_a(t)$ is $y_t=t(N-t)$.
In particular $2(B_a+D_a)=(m-1)(N+1)-x_a$.
The alternating completeness sum, obtained by substituting the
terminating series in \eqref{T}, is
\begin{equation}
 F(t):=\sum_r(-1)^rd_r\ell_r^2T_r(t)
   =(-1)^{m-1+t}mD_-(m-t).
\end{equation}
The same sum with $x_r$ inserted is
$[(m-1)(N+1)-y_t]F(t)$. Substitution in \eqref{mu} gives the
signed zeroth moment and the signed $x_r$ moment
$(x_a+\mathcal L_a)\xi_{ab}$.

For explicit reduction, the two charged components of a neutral
rank-$b$ tensor give
\begin{equation}
 \xi_{ab}=T_a(b)-\frac b{N+1-2b}[T_a(b-1)-T_a(b)].
\end{equation}
Thus $(\mathcal L_a-x_b)\xi_{ab}=
\frac{x_b}{N+1-2b}[T_a(b-1)-T_a(b)]$.
If $c_j$ is the $j$th summand of $\chi_{ab}$, the corresponding
summand of $\xi_{ab}$ is $mc_j/(m-j)$. Termwise subtraction shows
\begin{equation}
 \frac{x_b[T_a(b-1)-T_a(b)]}{N+1-2b}
 =-\sum_j\frac{jm}{m-j}c_j=-m(\xi_{ab}-\chi_{ab}),
\end{equation}
which proves the last formula in \eqref{moments}, including $b=0$.

\newpage
\section{Same-disorder cross and cancellation}
In the actual $p$-body ensemble let $d=m-p>0$, $C_a=\EE h_a^2$,
and $\omega_a=C_a/C$. Primitive embedding gives
\begin{equation}
 \sqrt{\frac{C_a^-}{C_a}}\sqrt{\frac{C_b^-}{C_b}}
       \frac{\ell_a\ell_b}{m}=\frac{d^2}{m}=:g.
\end{equation}
This independence of $a,b$ is useful. The upper half-filled product
obeys $P_r(V_bU_a)=\varepsilon P_r(U_aV_b)$.
Consequently the same-$J$ cross, including the $a=b$ Gaussian
decoupling coefficient, is
\begin{equation}
 \left\langle\sum f_{ab;r}\W Z^+_{ab;r},O\right\rangle_m
 =C^2\sum_{a,b}\omega_a\omega_b\sum_r f_{ab;r}(1+\varepsilon)
       [g\mu_{ab}(r)+(2p-m)\pi_{ab}(r)].\label{cross}
\end{equation}
Here $\pi$ is the previous upper-sector ordered-product probability;
pad $\mu$ by zero at $r=m$. The two terms in brackets must be
combined before an asymptotic expansion.

Set $\eta=p^2/m$, $\Delta=\xi_{ab}-\chi_{ab}$, and
$u_X=-2ab+x_ax_b/C_f$. The cross for $f=1$ and $f=X/N$,
per ordered input pair and divided by $C_aC_b$, is exactly
\begin{align}
 b_{ab}&=\eta(1+\chi_{ab})+g\Delta,\label{bexact}\\
 z_{ab}&=\frac1N\left\{
 \eta(u_X-2ab\chi_{ab})+
 g\left[\frac{x_ax_b}{m(m-1)(N+1)}+(m-2ab)\Delta\right]
 \right\}.\label{zexact}
\end{align}
We used $g+2p-m=\eta$ and
$C_t^{-1}-C_f^{-1}=[m(m-1)(N+1)]^{-1}$.

For fixed input deficits $p-a,p-b$ the terminating-series bounds
give
\begin{equation}
 \chi_{ab}\to q,\qquad
 m\Delta=\sum_j\frac{jm}{m-j}c_j\to-4\lambda q.
\end{equation}
Each $|c_j|$ is bounded by $K_\lambda^j/j!$, uniformly for
$0\leq a,b\leq p$ at sufficiently large $p$. Also $m/(m-j)$ is
uniformly bounded. Dominated summation therefore applies to the
series and its first $j$ moment. Equations \eqref{bexact}--\eqref{zexact}
then imply
\begin{equation}
 b_{ab}\to\theta(1-q),\qquad
 z_{ab}\to\theta^2(1+q)+\theta(1-q).
\end{equation}
They also give uniform bounds $|b_{ab}|+|z_{ab}|\leq K_\lambda$.
The previous factorial input-deficit tail passes these limits through
the full $\omega_a\omega_b$ sum. Multiplication by
$2\bar a-a-b$, with $\bar a=\sum a\omega_a$, gives the $L_N$ cross;
its limit vanishes by the same weighted-tail bound.

\newpage
\section{From the Casimir proxy to the actual deficit}
Define $\mathcal X_N=\sum X\W Z^+_{ab;r}$. The finite numerical
quantity in \eqref{zexact} is $\langle\mathcal X_N/N,O\rangle$.
To obtain \eqref{main}, use the exact identity
\begin{equation}
 X=S[N+1-2(a+b)+S].
\end{equation}
Orthogonality of input-pair and output-rank blocks, together with
the preceding positive parity Gram, gives
\begin{equation}
 \|\mathcal S_N-\mathcal X_N/N\|_{2,m}
 \leq\frac{4p}{N}\|\mathcal S_N\|_{2,m}=O(N^{-1/2}).\label{proxy}
\end{equation}
We need only a weak bound on the other factor. The previous
full-Fock, $\beta=\mu_{\rm chem}=0$ calculation bounds
$\|\W A\|_{2,\mathrm F}^2$ at fixed positive $\lambda$.
All sector contributions to this positive norm are nonnegative, so
\begin{equation}
 \|O\|_{2,m}^2\leq
 \frac{\|\W A\|_{2,\mathrm F}^2}{2^{-N}\binom N{N/2}}
 =O(N^{1/2}).\label{weakbound}
\end{equation}
Cauchy--Schwarz applied to \eqref{proxy} and \eqref{weakbound}
bounds the cross error by $O(N^{-1/4})\to0$.
This step does not identify the full-Fock norm with a fixed-sector
norm. The existing full-Fock result is used only as a positive bound.

\section{What the cross tells us}
Let $c=\theta(1-q)/(1+q)$, $\mathcal S_\perp=\mathcal S_N-cB$,
and $Y=O-cB$. The preceding parity Gram and \eqref{main} yield
\begin{align}
 \lim\langle\mathcal S_\perp,Y\rangle_m
 &=\lim\|\mathcal S_\perp\|_{2,m}^2=V_S,\notag\\
 V_S&=J_s^4\left[\frac{2(1-q)}\lambda+
                         \frac{16q}{1+q}\right]>0.\label{VS}
\end{align}
In particular $\langle\mathcal S_\perp,O-\mathcal S_N\rangle_m\to0$.
The output deficit is not orthogonal to the derivative singlet.
The remaining scalar quantity for this part of the projection is
\begin{equation}
 \Delta_N=\|O-\mathcal S_N\|_{2,m}^2.
\end{equation}
If it has a finite limit $\Delta_*$, the norm remaining after the
$B,O$ projection tends to $V_S\Delta_* /(V_S+\Delta_*)$.
This is a conditional projection formula, not an evaluation of
$\Delta_*$. The already controlled input direction is orthogonal
in this limiting Gram. Weighted $A$ tails, the fixed-sector $A$ norm,
and the low-chaos terms of the raw memory coefficient remain to be
computed. Only afterwards should the full energy-resolved spectrum
and bulk mass dictionary be assembled.

\newpage
\section{Verification and continuation}
All calculations here use the same CAR and covariance normalization
as the earlier memory notes. The validation includes:
\begin{itemize}
 \item 1,925 finite embedded/pure charged covariance identities at
 general charge, plus 139 orthogonality identities, 139 difference
 equations, and 556 signed/unsigned moments at half filling.
 \item Ten complete rectangular covariance maps, with 306,152
 complex entries, checked using exact CAR and Casimir projectors.
 These test the map, not merely selected highest-weight vectors.
 \item The full Gaussian quadratic-coefficient Gram on
 $(N,p,m)=(4,1,2),(6,1,3),(6,2,3)$: respectively 136, 666, and
 25,425 symmetric coordinate pairs. The $B$, $\mathcal S_N$, and
 $\mathcal X_N/N$ crosses with $O$ agree with \eqref{cross}.
 \item 384 contiguous first-moment identities and two symbolic
 rational identities; three complete input-rank sums compared with
 an independent exact pair summation.
 \item Twelve complete-rank double-scale calculations at
 $\lambda=0.1,0.25,0.5$ and $p=16,32,64,128$, evaluated at 80 and
 120 decimal digits. Every rank and every terminating-series term
 is retained; the cancellation factor is recorded.
\end{itemize}
The last calculation uses the Casimir proxy. For example at
$\lambda=0.5,p=128$, its cross divided by $C^2$ differs from its
limit by about $-0.0548285$; the $B$ and $L_N$ cross deviations
are about $1.95318\,10^{-5}$ and $3.32957\,10^{-5}$.
These are finite-$p$ deviations, not numerical errors. Precision
refinement is not an interval-rounding certificate. The proxy is
replaced by the actual deficit using the analytic bound on page 5.

The finite Gram gives positive residuals after projection onto
$B,O$ in all three test systems. This does not decide the
double-scaled residual. The next calculation should evaluate
$\Delta_N$ while preserving adjacent-charge interference, then
control the $g_A$ weighted tail from the previous core formula.
No new time-dependent spectral measure or new bulk particle is
claimed by this equal-time calculation.

\begin{thebibliography}{9}
\bibitem{kota} V. K. B. Kota and M. Vyas,
\emph{Random Matrix Theory for Transition Strength Densities in
Finite Quantum Systems: Results from Embedded Unitary Ensembles},
\href{https://arxiv.org/abs/1504.00759}{arXiv:1504.00759},
Secs. 3 and 7. The recoupling structure is used; its dilute limit
is not our half-filled double scaling.
\bibitem{ks} R. Koekoek and R. F. Swarttouw,
\emph{The Askey-scheme of hypergeometric orthogonal polynomials
and its $q$-analogue}, report 98-17 (1998), Sec. 1.2,
\href{https://tudelft.roelofkoekoek.nl/documents/as98.pdf}{Eq. (1.2.5)}.
\end{thebibliography}
\end{document}
